\documentclass[11pt]{article}
\usepackage[a4paper,margin=24mm]{geometry}
\usepackage{amsmath,amssymb,booktabs}
\usepackage[colorlinks=true,linkcolor=blue]{hyperref}
\title{Exact angular controls for nested six-dimensional $\phi^3$ graphs}
\author{Phi3Lab development record}
\date{25 October 2026}
\begin{document}
\maketitle
\noindent This document records algebraic variance reduction. The original forests,
local Taylor orders, signs and massive Euclidean conventions are unchanged.
An integrable majorant justifies rearrangement; it is not a numerical precision certificate.

\section{Full angular average for G371}
At fixed $t=q^2$ and $v=r^2$, write $h=q\cdot r$ and
\[
 a=1+t+v,\qquad b=2\sqrt{tv},\qquad h=\frac b2\xi,\qquad
 d\mu_6(\xi)=\frac{8}{3\pi}(1-\xi^2)^{3/2}\,d\xi .
\]
The perpendicular four-dimensional relative angle used by the complete density
has semicircle measure $d\mu_4=(2/\pi)\sqrt{1-\xi^2}\,d\xi$.
Thus $d\mu_6=\frac43(1-\xi^2)d\mu_4$.
Define $R=\sqrt{(a-b)(a+b)}$ and $S=a+R$. The inverse moments are
\begin{align*}
 L_1&=\frac{4(a+2R)}{3S^2},&
 L_2&=\frac4{S^2},&
 L_3&=\frac4{RS^2},&
 L_4&=\frac{4(a+2R)}{3R^3S^2}.
\end{align*}
They satisfy $L_{n+1}=-(\partial_a L_n)/n$. The mass protects the endpoint:
$a-b=1+(\sqrt t-\sqrt v)^2>0$.

Let the four non-cross host ratios times the two complete triangle factors be
$F_0+\lambda F_1+\lambda^2F_2+\cdots$, with $P_0=F_0$.
The two cross denominators contribute
$[X(X+\lambda\ell+\lambda^2)]^{-1}$, where
$X=a+2h$ and $\ell=2(q+r)\cdot(p_1+p_2)$. The quadratic coefficient is
\[
 P_2=\frac{F_2}{X^2}-\frac{\ell F_1+P_0}{X^3}
       +\frac{\ell^2P_0}{X^4}.
\]
Rotation at fixed $(t,v,h)$ gives
\[
 \langle F_2\rangle=A_F+B_Fh,\quad
 \langle\ell F_1\rangle=A_L+B_Lh,\quad
 \langle\ell^2\rangle=\frac23(t+v+2h).
\]
Put $H_n=(L_{n-1}-aL_n)/2$. The complete angular mean is
\[
 G=A_FL_2+B_FH_2-A_LL_3-B_LH_3-P_0L_3
       +\frac23P_0(L_3-L_4).
\]
Let $\mathcal F$ be the complete perpendicular-angle marginal, and
$\mathcal P$ its pointwise quadratic marginal. If
$w=(1-z_1)(1-z_2)$ and $M$ is the radial density, the new density is
\[
 \frac{M}{(1+t)^2(1+v)^2}
       \left[w(\mathcal F-\mathcal P)+\frac G4\right].
\]
The factor $1/4$ is exact because $\int_0^1\int_0^1w\,dz_1dz_2=1/4$.
This avoids sampling an already known angular average.

\section{Fourth-order control for G203 and G64}
Keep both local two-point $T^2$ operations and all vertex $T^0$ operations.
The fourth-order coefficient below is only a zero-integral angular control.
After analytic child integration, the physical product is
\[
 H(\lambda)K(V_0+\lambda L_p+\lambda^2C_p)
 J(t+\lambda L_b+\lambda^2C_b),\qquad J(t)=\frac{B(t)}{1+t},
\]
where the extra denominator in $J$ is the duplicated host propagator.
The bubble and parent derivative are
\[
 B(t)=\int_0^1[(1+tu(1-u))\log(1+tu(1-u))-tu(1-u)]\,du,
 \qquad K'(V)=-\frac{x B((1-x+V)/x)}{(1+V)^2}.
\]
The coefficient of $\lambda^4$, denoted $P_4$, is formed by truncated power
series, retaining every product term. No finite difference or new local
subtraction is introduced.

For a planar polynomial $f$ of degree at most four, the SO(6) mean at radius
$\sqrt t$ is exactly
\[
 \langle f\rangle=\frac13f(0,0)+\frac1{12}
  \sum_{k=0}^{7} f\left(\sqrt{t/2}\cos\frac{k\pi}{4},
                         \sqrt{t/2}\sin\frac{k\pi}{4}\right).
\]
The independent reference uses twelve axis directions of weight $-1/48$
and sixty signed pair directions of weight $1/48$ in six dimensions.
Both rules reproduce all fifteen monomials through total degree four:
\[
 \langle q_x^2\rangle=\frac t6,\qquad
 \langle q_x^4\rangle=\frac{t^2}{16},\qquad
 \langle q_x^2q_y^2\rangle=\frac{t^2}{48}.
\]
The added density is
$\mathrm{common}\,[\langle P_4\rangle/2-(1-z)P_4]$.
Its angular integral vanishes exactly.

\section{An explicit integrability bound}
For $t\ge16$, let $L=\log(1+t)$, $d=1+V_0$ and $U=\log(d/x)$.
The simplex representation yields
\[
 |B^{(j)}(t)|\le\frac{(j-2)!}{6t^{j-1}}\quad(j\ge2),\qquad
 B\le\frac{tL}{6},\quad B'\le\frac L6.
\]
Consequently the $\lambda^n$ coefficient of $J$ is bounded by
$3^nL/t^{n/2}$ for $0\le n\le4$, and the host coefficient is bounded by
$\binom{n+2}{2}3^n/t^{n/2}$.
For the parent,
\[
 \frac{|K^{(n)}(V_0)|}{n!}\le\frac{1+U}{d^n}\quad(1\le n\le4).
\]
The endpoint-safe inequality is
$x+(d-x)u(1-u)\ge d\,u(1-u)$.
Together with $|L_p|^2\le4V_0C_p$ and $C_p\le1/3$, this bounds its
positive-order $\lambda^j$ coefficient by $3^j(1+U)/d^{j/2}$.

For these two routed graphs, $V_0=t\xi(1-\xi)$ with $\xi=(1-x)y$.
The two parameter moments needed are $\pi/2<2$ and $\pi\log2<3$.
Thus a positive-order parent coefficient integrates to at most
$8\,3^jL/\sqrt t$, while the zeroth-order term integrates to at most $L^2/6$.
Summing the fourth-order product gives
\[
 \int |P_4|\,d\text{(parent parameters)}
 \le \frac{22680L^2}{\sqrt t}+\frac{945L^3}{2t^2}
 <\frac{23153L^3}{\sqrt t}.
\]
After the host radial measure, mean-minus-point and the previous control,
a valid whole-tail majorant is $50000(1+\log t)^3/t^{3/2}$.
For $T\ge16$, its integrated tail is
\[
 \frac{100000}{\sqrt T}
 [\,\ell^3+6\ell^2+24\ell+48\,],\qquad \ell=1+\log T.
\]
Massive continuity and the integrable $\log x$ endpoint cover finite radii.

\section{Evidence and acceptance}
Proofs rebuild the original forests, verify exact identities and retain
all source hashes. Native ball arithmetic is compared with independently
differentiated scalar references and direct auxiliary quadratures.
Every production request must be replayed before accounting.
The empirical acceptance gates and sample allocation remain fixed.
The actual job receipts, not the existence of this document, determine status.
\end{document}